\section*{Bab \ref{ch:g5:simple-machines} --- Pesawat Sederhana: Tuas, Katrol, Bidang Miring}

\begin{solution}{exo:g5:simple-machines:1}
\omterm{def:g3:levers-and-scales:lever}{Tuas} (linggis, jungkat-jungkit, gunting), \omterm{def:g5:simple-machines:pulley}{katrol} (tiang bendera,
sumur, derek), dan \omterm{def:g5:simple-machines:ramp}{bidang miring} (landasan kursi roda, jalan gunung,
landasan bongkar muat).
\end{solution}

\begin{solution}{exo:g5:simple-machines:2}
\omterm{def:g5:simple-machines:pulley}{Katrol} tetap hanya mengubah arah tarikannya (tarik ke bawah, bebannya
naik). \omterm{def:g5:simple-machines:pulley}{Katrol} bergerak mengubah besarnya: bebannya tergantung pada dua
untai, sehingga tangannya hanya menahan separuh \omterm{def:g4:weight-and-mass:weight}{beratnya}.
\end{solution}

\begin{solution}{exo:g5:simple-machines:3}
Menarik ke bawah mengerahkan beratmu sendiri: kamu dapat bersandar
atau bahkan bergantung pada talinya, dan membiarkan tarikan Bumi
menolong ototmu. Kalau mengangkat ke atas, lenganmu berjuang sendiri.
\end{solution}

\begin{solution}{exo:g5:simple-machines:4}
Separuh \omterm{def:g4:weight-and-mass:weight}{beratnya}. Untuk menaikkan bebannya \qty{3}{m}, kedua untainya
harus memendek \qty{3}{m}: tanganmu menarik $2 \times 3 = \qty{6}{m}$
tali.
\end{solution}

\begin{solution}{exo:g5:simple-machines:5}
Landasan yang curam meregangkannya lebih jauh. Tenaga yang lebih kecil
pada landasan yang landai dibayar dengan jarak: tinggi yang sama
dicapai di sepanjang jalan yang jauh lebih panjang --- itulah
pertukaran aturan emas.
\end{solution}

\begin{solution}{exo:g5:simple-machines:6}
Lurus mendaki lerengnya, tarikan yang diperlukan akan mengalahkan
mesin atau pasukan bagal mana pun; kelokannya melilitkan \omterm{def:g5:simple-machines:ramp}{bidang miring}
yang panjang dan landai menyeberangi lereng gunung --- tenaga kecil,
jalan panjang. Seluruh jalannya adalah satu landasan miring yang
direntangkan.
\end{solution}

\begin{solution}{exo:g5:simple-machines:7}
Tidak ada pesawat yang memberi sesuatu dengan cuma-cuma: \omterm{def:g2:pushes-and-pulls:force}{gaya} yang
dihemat dibayar kembali dalam jarak, dengan persis. Sepersepuluh \omterm{def:g2:pushes-and-pulls:force}{gaya}
memakan sepuluh kali jalannya (atau talinya).
\end{solution}

\begin{solution}{exo:g5:simple-machines:8}
Ulirnya adalah lereng panjang yang landai yang dililitkan pada
batangnya. Banyaknya putaran obeng yang ringan itulah jalan yang
panjang; turunnya sekrup yang mungil pada tiap putaran adalah tinggi
kecil yang diperoleh.
\end{solution}

\begin{solution}{exo:g5:simple-machines:9}
Seperempat \omterm{def:g4:weight-and-mass:weight}{beratnya}. Masing-masing dari keempat untainya harus
memendek \qty{2}{m}: $4 \times 2 = \qty{8}{m}$ tali melewati tangan
pekerjanya.
\end{solution}

\begin{solution}{exo:g5:simple-machines:10}
Landasan yang baru menaiki \qty{1}{m} yang sama dalam separuh jalan,
jadi ia dua kali lebih curam, dan dorongannya kira-kira menjadi dua
kali lipat. Aturan emas menagih seketika: jalannya diparuh, \omterm{def:g2:pushes-and-pulls:force}{gayanya}
dua kali lipat --- pendakian petinya memakan biaya yang sama dengan
cara mana pun, dan para pengangkut sekadar menukar banyak uang receh
yang nyaman dengan beberapa keping yang \omterm{def:g4:weight-and-mass:weight}{berat}.
\end{solution}

\begin{solution}{exo:g5:simple-machines:11}
Tidak ada susunan yang dapat memenuhinya. Menaikkan bebannya satu
\omterm{def:g2:measuring-length:units}{meter} memakan jatah \omterm{def:g5:everyday-energy:energy}{energi} yang tetap; kalau dibayar dengan separuh
\omterm{def:g2:pushes-and-pulls:force}{gaya}, penjualnya harus menarik \emph{dua} \omterm{def:g2:measuring-length:units}{meter} tali --- separuh \omterm{def:g2:pushes-and-pulls:force}{gaya}
di sepanjang satu \omterm{def:g2:measuring-length:units}{meter} saja hanya menyerahkan separuh tagihannya.
Iklan itu menjanjikan angkatan berbayar dengan harga separuh: persis
makan siang gratis yang dilarang aturan emas, dan yang dahulu
menggagalkan mesin abadi.
\end{solution}
